% =====================================================================
% §2.5 — La molteplicità delle voci (esempi: ex4_voices, ex5_scatter)
%
% Pattern: risposta al ponte di §2.4 -> quattro assi ortogonali +
%          strategie (DISCIPLINA: nominarne due, registry una volta) ->
%          ex4 = polo deterministico (bit-identico) -> scatter come
%          accoppiamento temporale -> ex5 -> ponte a §2.6
%
% TODO prima della submission:
%   [ ] ex4_voices.yml: fissare num_voices a 5 costante (l'attuale envelope
%       [[0,5],[.5,1],[1,10]] era sperimentale: per la figura a cinque bande
%       parallele e la lettura piu' pulita serve N fisso) e duration 20
%   [ ] linerange dei listati: allineare ai file reali (ex4: blocco voices;
%       ex5: le sole righe distribution + scatter)
%   [ ] verificare nomi figure: ex4_voices_map.pdf, ex5_scatter_map.pdf
%   [ ] B&W check: la separazione delle bande su Y (pointer linear) deve
%       reggere la stampa in scala di grigi anche senza il colore=pitch
%   [ ] valvola di compressione (se la sezione sfora): fondere le due
%       figure in un doppio pannello o ridurre scatter a sola prosa
% =====================================================================

% ----------------------------------------------------------------
\subsection{La molteplicità delle voci}\label{sec:voci}

Moltiplicare la voce significa dichiarare che lo stream non emette un
grano per volta ma $N$ grani coordinati. Se le $N$ voci fossero
identiche, la moltiplicazione sarebbe un puro guadagno d'energia; la
domanda compositiva è come differenziarle senza doverle scrivere una per
una.

Il blocco \texttt{voices} risponde con quattro assi ortogonali di
differenziazione --- trasposizione, ritardo d'attacco, posizione di
lettura, posizione spaziale --- ciascuno governato da una strategia
intercambiabile (il registro è estensibile); la voce~0 è il riferimento
e non riceve offset, le altre si dispongono rispetto a essa.
L'esempio (Listato~\ref{lst:voci}) ne impiega due: una strategia
d'accordo sulla trasposizione e una lineare su lettura e spazio.

    \lstinputlisting[
      linerange={30-41}, % TODO: allineare al file reale
      language=yaml,
      basicstyle=\ttfamily\fontsize{7.5}{9}\selectfont,
      label=lst:voci,
      caption={Il blocco \texttt{voices}: cinque voci differenziate su
      tre assi (\texttt{ex4\_voices.yml}).}
    ]{examples/ex4_voices/ex4_voices.yml}

Cinque voci ricevono i cinque gradi di un accordo di
nona\footnote{L'accordo è qui specificato in semitoni; il sistema
ammette unità microtonali (cents, EDO arbitrari).} --- nella partitura,
dove il colore codifica la trasposizione, cinque colori distinti ---
cinque teste di lettura equidistanti, che sull'asse~Y della
Fig.~\ref{fig:voci-score} si separano in altrettante bande parallele, e
un ventaglio spaziale di 150\textdegree, udibile ma estraneo agli assi
della partitura. Si noti il regime: dopo due sottosezioni di
stocasticità governata, qui ogni strategia è deterministica e il
rendering è bit-identico fra le esecuzioni.

\begin{figure}[h]
    \includegraphics[width=\linewidth]{examples/ex4_voices/ex4_voices_map}\\
    \captionof{figure}{Partitura di \texttt{ex4\_voices}: cinque bande
    parallele sull'asse della posizione di lettura, una per voce; il
    colore codifica la trasposizione (i cinque gradi dell'accordo).}
    \label{fig:voci-score}
\end{figure}

Le voci così ottenute condividono però ancora la stessa griglia
temporale: i loro grani nascono insieme, sugli stessi istanti della voce
di riferimento. La moltiplicazione non è ancora indipendenza, e
l'indipendenza ha nel sistema un proprio asse: \texttt{scatter}, il
grado di accoppiamento temporale fra le voci. A 0 tutte le voci
condividono la sequenza di intervalli della voce~0; a 1 ciascuna
calcola la propria; i valori intermedi interpolano. Nell'esempio
(Listato~\ref{lst:scatter}) l'unica variabile è l'envelope su
\texttt{scatter}; la distribuzione temporale è tenuta asincrona e
\emph{costante} --- condizione di servizio dichiarata: scatter agisce
solo se la griglia ha contenuto stocastico da diversificare
(\S\ref{sec:griglia}) --- e la densità è bassa perché i singoli onset
restino leggibili.

    \lstinputlisting[
      linerange={36-38}, % TODO: allineare al file reale (distribution + scatter)
      language=yaml,
      basicstyle=\ttfamily\fontsize{7.5}{9}\selectfont,
      label=lst:scatter,
      caption={Le righe decisive di \texttt{ex5\_scatter.yml}: griglia
      asincrona costante, envelope sul solo \texttt{scatter}.}
    ]{examples/ex5_scatter/ex5_scatter.yml}

Nella Fig.~\ref{fig:scatter-score} l'esito è leggibile sull'allineamento
degli onset: a sinistra, con le griglie agganciate, i grani delle
quattro bande cadono sulle stesse ascisse --- colonne verticali che
attraversano le voci; procedendo verso destra le colonne si sfaldano e
ogni banda segue il proprio cammino temporale. All'ascolto, da un
ensemble che articola insieme a quattro flussi che si sfilano l'uno
dall'altro. Qui il non-determinismo rientra, e con esso la
riproducibilità per andamento: le voci si sganciano a ogni esecuzione,
mai negli stessi istanti.

\begin{figure}[h]
    \includegraphics[width=\linewidth]{examples/ex5_scatter/ex5_scatter_map}\\
    \captionof{figure}{Partitura di \texttt{ex5\_scatter}: gli onset
    delle quattro voci, allineati in colonne a sinistra
    (\texttt{scatter}=0), si disaccoppiano verso destra
    (\texttt{scatter}=1).}
    \label{fig:scatter-score}
\end{figure}

Lo stream, le sue deviazioni e le sue voci sono ora interamente
dichiarati. Resta il percorso inverso: come la dichiarazione diventa
suono --- e come il sistema tiene questo percorso abbastanza corto da
poterlo percorrere molte volte.



